\documentclass[11pt]{article}
\usepackage[margin=0.8in]{geometry}
\usepackage[T1]{fontenc}
\usepackage{lmodern,amsmath,amssymb,amsthm,booktabs,array}
\usepackage[colorlinks=true,linkcolor=blue,urlcolor=blue,citecolor=blue]{hyperref}
\hypersetup{pdftitle={Five conjecture resolutions in three families},pdfauthor={Blake Stone},pdfsubject={Polynomial irreducibility, factorial-ratio divisibility, replacement words, and complete Lean proofs}}
\setlength{\parindent}{0pt}
\setlength{\parskip}{0.38em}
\newtheorem{theorem}{Theorem}[section]
\newtheorem{lemma}[theorem]{Lemma}
\newcommand{\N}{\mathbb N}
\newcommand{\Z}{\mathbb Z}
\newcommand{\Q}{\mathbb Q}
\newcommand{\R}{\mathbb R}
\begin{document}
\begin{center}
{\LARGE\bfseries Five conjecture resolutions\\[3pt]in three families}\par
\vspace{0.5em}
{\large Blake Stone}\par
Proof developed with ChatGPT\par
November 27, 2026
\end{center}

\textbf{Abstract.}
We prove the full irreducibility classification for the normalized
derivative numerators of $1/(1-x+x^2)$: the nonconstant polynomial of
degree $n$ is irreducible over $\Q$ exactly when $n+1$ is a prime other
than three. A reciprocal Eisenstein argument handles every allowed
prime; classical Lucas divisibility excludes composite indices.
We also prove the entire positive-shift integrality family in
A211417, giving the explicit uniform multiplier $((30r)!)^r$ for
each $k\in\{2,3,5\}$ and every shift range $r$.
Finally, an exact description of a binary replacement word proves
its conjectured length recurrence and refutes both proposed
occurrence-position limits. The actual zero-position ratios have
subsequential limits $5/2$ and $7/3$, and the one-position ratios
have subsequential limits $5/3$ and $7/4$.
Three Lean projects verify all five assertions from the literal
source definitions, using fifteen local proof modules and pinned
mathlib dependencies. Known ingredients and earlier neighboring
results are credited. Bounded prior-publication checks found no
earlier proof of the exact assertions claimed here; this is not a
certification of historical priority or editorial acceptance.

\begin{center}
\begin{tabular}{@{}l>{\raggedright\arraybackslash}p{4.55in}@{}}
\toprule
Source & Result and count\\
\midrule
A326926/A326933 & One full irreducibility classification, repeated in two entries.\\
A211417 & One general positive-shift integrality conjecture, covering all three prescribed values of $k$ and every $r$.\\
A289128/A289131 & One shared conjecture identifying the literal iterate lengths with an independently specified recurrence.\\
A289129 & One limit refutation, strengthened to nonexistence of any real limit.\\
A289130 & One limit refutation, strengthened to nonexistence of any real limit.\\
\bottomrule
\end{tabular}
\end{center}

\section{Prime-index irreducibility of the derivative polynomials}
\label{sec:polynomials}

Kimberling's A326926~\cite{A326926} takes the numerator of the
derivative of order $n$ of $1/(1-x+x^2)$ and divides it by $n!$.
The row has degree $n$, beginning with the constant row at $n=0$.
The irreducibility comment is repeated in A326933~\cite{A326933};
these are the same assertion. Peter Bala's November 2026 addendum
already records the recurrence, closed form, roots, and strong
divisibility. Its one-based row $R(N,x)$ is our $P_{N-1}(x)$.

Put
\[
D=1-x+x^2,\qquad A=1-2x,\qquad
P_0=1,\quad P_1=A,\quad P_{n+2}=AP_{n+1}-DP_n.
\]

\begin{theorem}\label{thm:polynomials}
For every $n\ge0$,
\[
\left(\frac1D\right)^{(n)}=\frac{n!P_n}{D^{n+1}},
\qquad \gcd(P_n,D)=1,\qquad \deg P_n=n.
\]
The canonical reduced numerator over $\Q[x]$ is exactly $n!P_n$.
For every $n>0$,
\[
P_n\text{ is irreducible over }\Q
\quad\Longleftrightarrow\quad n+1\text{ is prime and }n+1\ne3.
\]
\end{theorem}

\subsection{Differential identities and exact reduced numerators}

The recurrence implies the two identities
\begin{align}
P_{n+1}'&=-(n+2)P_n,\label{eq:appell}\\
DP_n'+(n+1)AP_n&=(n+1)P_{n+1}.\label{eq:differential}
\end{align}
For a simultaneous induction, both hold at $n=0$. Since $A'=-2$
and $D'=-A$, differentiation of the recurrence and the two
inductive hypotheses give
\[
\begin{aligned}
P_{n+2}'
 &=-2P_{n+1}+AP_{n+1}'+AP_n-DP_n'\\
 &=-2P_{n+1}-(n+1)AP_n-DP_n'
 =-(n+3)P_{n+1}.
\end{aligned}
\]
Also
\[
DP_{n+1}'+(n+2)AP_{n+1}
=(n+2)(AP_{n+1}-DP_n)=(n+2)P_{n+2}.
\]
This proves both identities for every $n$.

The derivative formula in Theorem~\ref{thm:polynomials} follows
by induction from the quotient rule:
\[
\left(\frac{n!P_n}{D^{n+1}}\right)'
=\frac{n!\bigl(DP_n'+(n+1)AP_n\bigr)}{D^{n+2}}
=\frac{(n+1)!P_{n+1}}{D^{n+2}}.
\]
The initial expression is $1/D$.

To verify that this is the reduced numerator, observe
\[
A^2=4D-3,\qquad -\frac A3 A+\frac43D=1.
\]
Thus $A$ and $D$ are coprime over $\Q$. The base rows $1,A$
are coprime to $D$, and the recurrence gives
$P_{n+2}\equiv AP_{n+1}\pmod D$. Induction proves every $P_n$
coprime to $D$. The scalar $n!$ is a rational unit, so $n!P_n$
is coprime to $D^{n+1}$. This denominator is monic, which fixes
the scale of the canonical numerator to be exactly $n!P_n$.
The source's normalization by $n!$ therefore recovers precisely
the recurrence polynomial.

\subsection{Coefficients needed for Eisenstein}

The recurrence gives $\deg P_n\le n$ and
\begin{equation}\label{eq:leading}
[x^n]P_n=(-1)^n(n+1).
\end{equation}
Indeed, the top coefficient in the next recurrence step is
\[
(-2)(-1)^{n+1}(n+2)-(-1)^n(n+1)
=(-1)^{n+2}(n+3).
\]
It is nonzero in characteristic zero, proving the degree equality.

Let $s_0=0,s_1=1,s_{j+2}=s_{j+1}-s_j$. Specialization at zero gives
$P_n(0)=s_{n+1}$. Six applications of this recurrence give
$s_{j+6}=s_j$ and
\[
(s_0,s_1,s_2,s_3,s_4,s_5)=(0,1,1,0,-1,-1).
\]
For a prime $p\ne3$, either $p=2$ or $p\equiv1,5\pmod6$.
Consequently $P_{p-1}(0)=\pm1$.

Now write $P_{p-1}=\sum_{j=0}^{p-1}c_jx^j$. The derivative identity
\eqref{eq:appell} says $P_{p-1}'=-pP_{p-2}$. Coefficient comparison
shows $p\mid jc_j$ whenever $1\le j\le p-1$. Since $p\nmid j$,
we have $p\mid c_j$ throughout that range. Its final coefficient
is $c_{p-1}=(-1)^{p-1}p$, so $p^2\nmid c_{p-1}$.

Thus the reciprocal polynomial
\[
Q(x)=x^{p-1}P_{p-1}(x^{-1})
=c_0x^{p-1}+c_1x^{p-2}+\cdots+c_{p-1}
\]
is Eisenstein at $p$: its leading coefficient is $\pm1$, all other
coefficients are divisible by $p$, and its constant coefficient
is not divisible by $p^2$. The classical Eisenstein criterion and
Gauss's lemma prove $Q$ irreducible over $\Q$. A factorization of
$P_{p-1}$ would have factors with nonzero constant coefficients;
their reciprocals would factor $Q$ with unchanged degrees.
Therefore $P_{p-1}$ is irreducible.

\subsection{Composite indices and the exceptional prime}

Let $U_0=0,U_1=1,U_{j+2}=AU_{j+1}-DU_j$. Then $U_N=P_{N-1}$
for $N\ge1$. The classical Lucas addition identity is
\[
U_{m+n+1}=U_{m+1}U_{n+1}-DU_mU_n.
\]
It follows directly by induction on either index. In particular,
\[
U_{kd+d}=U_{kd}U_{d+1}-DU_{kd-1}U_d
\]
proves by induction on $k\ge1$ that $U_d\mid U_{kd}$. Hence
\[
d\mid N\quad\Longrightarrow\quad P_{d-1}\mid P_{N-1}
\qquad(d,N>0).
\]
These identities are existing results, including Norfleet's
Theorem 1 and Corollary 1~\cite{Norfleet}; their formal proofs
are included as dependencies.

If $N=n+1$ is composite, it has a proper prime divisor $d$.
The divisor $P_{d-1}$ then has degree strictly between zero and
$N-1$. It is a nonconstant proper factor of $P_n$, proving
reducibility. Finally $P_2=3x(x-1)$ deals with $N=3$.
This completes Theorem~\ref{thm:polynomials}.

\paragraph{Formal endpoints.}
The polynomial classification is
\path{SixthDerivative.irreducible_iff}. The source bridge proves
\path{derivative_numerator} and
\path{normalized_source_numerator}, including coprimality and monic
normalization. The final endpoints
\path{SixthDerivative.source_numerator_irreducible_iff} and
\path{SixthDerivative.source_irreducible_iff} state the criterion
for the canonical numerator and the exact normalized source row.
The representation-independent rational-function derivative module
is reused unchanged from the preceding package. No assertion is
made here about the exact number of factors in every reducible row.

\section{An explicit multiplier for every positive shift range}
\label{sec:factorial}

The classical integral factorial ratio A211417 is
\begin{equation}\label{eq:factorial-ratio}
a(n)=\frac{(30n)!\,n!}{(15n)!\,(10n)!\,(6n)!},\qquad n\ge0.
\end{equation}
Peter Bala's August 28, 2025 conjecture asks, for each
$k\in\{2,3,5\}$ and $r\ge1$, whether a positive constant
$C(k,r)$ makes
\[
\frac{C(k,r)a(n)}{
\displaystyle\prod_{\substack{1\le i\le r\\\gcd(i,k)=1}}(kn+i)}
\]
an integer for every $n\ge0$~\cite{A211417}. We prove that the
same explicit choice works for all three values of $k$.

\begin{theorem}\label{thm:factorial}
For $k\in\{2,3,5\}$ and all $r,n\ge0$, put
\[
S=\{i:1\le i\le r,\ \gcd(i,k)=1\},\qquad
B=\prod_{i\in S}(kn+i).
\]
Then
\[
\boxed{B\mid ((30r)!)^r a(n).}
\]
In particular, $C(k,r)=((30r)!)^r>0$ proves the entire
positive-shift existence conjecture. The case $r=0$ uses an empty
product and is also included in the formal theorem.
\end{theorem}

\subsection{The nonnegative floor contribution}

For an integer $m\ge1$, define
\[
F(n,m)=\left\lfloor\frac{30n}{m}\right\rfloor
+\left\lfloor\frac n m\right\rfloor
-\left\lfloor\frac{15n}{m}\right\rfloor
-\left\lfloor\frac{10n}{m}\right\rfloor
-\left\lfloor\frac{6n}{m}\right\rfloor.
\]
If $t=\lfloor30n/m\rfloor$, nested integer division gives
\[
F(n,m)=H(t),\qquad
H(t)=t+\left\lfloor\frac t{30}\right\rfloor
-\left\lfloor\frac t2\right\rfloor
-\left\lfloor\frac t3\right\rfloor
-\left\lfloor\frac t5\right\rfloor.
\]
For example $\lfloor t/2\rfloor=\lfloor15n/m\rfloor$;
the other terms follow in the same way. We have
$H(t+30)=H(t)$. Its thirty residues are displayed below:
\[
\begin{array}{c|rrrrrrrrrr}
t=10q+s&s=0&1&2&3&4&5&6&7&8&9\\\hline
q=0&0&1&1&1&1&1&0&1&1&1\\
q=1&0&1&0&1&1&0&0&1&0&1\\
q=2&0&0&0&1&0&0&0&0&0&1
\end{array}
\]
Thus $F(n,m)\in\{0,1\}$, in particular $F(n,m)\ge0$.
Legendre's factorial formula gives
\begin{equation}\label{eq:legendre-a}
v_p(a(n))=\sum_{e\ge1}F(n,p^e).
\end{equation}
More explicitly, the nonnegative difference of the numerator and
denominator valuations proves the known integrality of
\eqref{eq:factorial-ratio} first; equation~\eqref{eq:legendre-a}
then follows for its integer value. Every sum is finite.

\subsection{A large divisor forces a unit contribution}

Let $m>30r$. At most one $i\in S$ can satisfy $m\mid kn+i$:
two distinct such indices would give $m\mid i-j$ despite
$0<|i-j|<m$.

Suppose such an $i$ exists and write $kn+i=mj$. Put $d=30/k$.
Since $\gcd(i,k)=1$ and $k>1$, $k\nmid i$; if $k\mid j$,
the equation would imply $k\mid i$. Hence $k\nmid j$.
Multiplying the equation by $d$ gives
\[
30n+di=mdj,\qquad 0<di<m.
\]
It follows that $t=\lfloor30n/m\rfloor=dj-1$. The possible
residues of $t$ modulo thirty are therefore
\[
\begin{array}{c|c|c}
k&j\bmod k&t\bmod30\\\hline
2&1&14\\
3&1,2&9,19\\
5&1,2,3,4&5,11,17,23
\end{array}
\]
Every displayed residue has $H(t)=1$. Thus the presence of a
large divisor of any requested linear factor is paid for by a
contribution $F(n,m)=1$ in the factorial-ratio valuation.

\subsection{All factors, including multiplicities}

Write
\[
N_m=\#\{i\in S:m\mid kn+i\}.
\]
For $m>30r$, the preceding argument proves $N_m\le F(n,m)$.
For $1\le m\le30r$, we have $N_m\le|S|\le r$ and
$\lfloor30r/m\rfloor\ge1$. Since $F(n,m)\ge0$, both cases give
the single bound
\begin{equation}\label{eq:contribution}
N_m\le F(n,m)+r\left\lfloor\frac{30r}{m}\right\rfloor.
\end{equation}
For any prime $p$, counting every power dividing every factor yields
\[
v_p(B)=\sum_{e\ge1}N_{p^e}.
\]
Summing~\eqref{eq:contribution} and using Legendre's formula gives
\[
\begin{aligned}
v_p(B)
&\le\sum_{e\ge1}F(n,p^e)
+r\sum_{e\ge1}\left\lfloor\frac{30r}{p^e}\right\rfloor\\
&=v_p(a(n))+r\,v_p((30r)!)
=v_p\bigl(((30r)!)^r a(n)\bigr).
\end{aligned}
\]
The valuation inequality for every prime proves the divisibility
in Theorem~\ref{thm:factorial}. All factors $kn+i$ are positive,
so the literal rational quotient is defined and integral. This
argument accounts for the full product and repeated prime factors.

\paragraph{Formal scope and prior work.}
The Lean namespace \path{ChebyshevShifts} proves the contribution
bound, the finite valuation sums, and
\path{positive_shift_divisibility}. It defines the source ratio
literally in $\Q$ and proves \path{sourceRatio_eq_a}, so natural
division is not being used as a substitute for integrality.
The theorem \path{positive_shift_rational_integrality} proves the
stated quotient is the image of an integer. The existential source
endpoint is \path{bala_positive_shift_conjecture}.

The factorial ratio's original integrality is classical and is not
counted as a new result. Earlier proofs of four fixed-factor
assertions and the separate negative-shift family are also excluded
from the contribution. The explicit multiplier here is generous;
we do not claim it is the least possible multiplier. The source and
prior-work report in the package records the checked earlier proofs
and the exact scope of the positive-shift claim.

\section{A replacement word whose position ratios do not converge}
\label{sec:a289128}

Clark Kimberling's entries A289128--A289131~\cite{A289128,A289129,A289130,A289131}, introduced on June 28,
2017, concern the left-to-right nonoverlapping replacement
\[
00\longmapsto0010,\qquad01\longmapsto011,\qquad10\longmapsto010.
\]
Let $R$ denote one pass of this scanner. A matched pair is consumed,
and newly emitted letters are not scanned during that pass. An
unmatched input letter is copied; in particular, an initial $11$
causes the first $1$ to be copied before scanning continues. Set
$w_0=00$ and $w_{n+1}=R(w_n)$. The first words are
\[
00,\quad0010,\quad0010010,\quad00100100110,
\]
followed by $001001001100101010$ and
$0010010011001010100110110110$.

The entry A289131 conjectures that the lengths satisfy its printed
order-five recurrence. The adjacent entries A289129 and A289130
conjecture that the positions of the $k$-th zero and one, divided by
$k$, tend respectively to $5/2$ and $5/3$. We prove the length
assertion and prove that neither position ratio has any real limit.
The repeated length assertion in A289128 is counted once.

\paragraph{Theorem.}
The literal words $w_n$ are nested prefixes and determine a unique
infinite binary word $W$. Write $L_n=|w_n|$ and let $Z_n,O_n$ be the
counts of its zeros and ones. For every $m\ge0$,
\[
\begin{array}{c|cc}
 & n=2m+1 & n=2m+2\\ \hline
L_n & 10\cdot2^m-3m-6 & 14\cdot2^m-3m-7\\
Z_n & 4\cdot2^m-1 & 6\cdot2^m-1\\
O_n & 6\cdot2^m-3m-5 & 8\cdot2^m-3m-6
\end{array}
\tag{\ref{sec:a289128}.1}
\]
The exceptional initial length is $L_0=2$. With the first six values
$2,4,7,11,18,28$, these lengths satisfy exactly
\[
L_n=L_{n-1}+3L_{n-2}-3L_{n-3}-2L_{n-4}+2L_{n-5}
\qquad(n\ge6).
\tag{\ref{sec:a289128}.2}
\]
Let $a(k)$ and $b(k)$ be the one-based positions of the $k$-th zero
and one in $W$. Both are defined for every $k\ge1$. Their ratios have
the subsequential limits
\[
\begin{aligned}
\frac{a(Z_{2m+1})}{Z_{2m+1}}&\longrightarrow\frac52,
&\frac{a(Z_{2m+2})}{Z_{2m+2}}&\longrightarrow\frac73,\\
\frac{b(O_{2m+1})}{O_{2m+1}}&\longrightarrow\frac53,
&\frac{b(O_{2m+2})}{O_{2m+2}}&\longrightarrow\frac74.
\end{aligned}
\tag{\ref{sec:a289128}.3}
\]
All four occurrence-rank subsequences are strictly increasing.
Consequently neither $a(k)/k$ nor $b(k)/k$ converges.

\paragraph{Proof: exact blocks and scanner boundaries.}
For $t\ge0$, define $A_t=0(10)^t$ and $B_t=0(110)^t$. For any finite
binary suffix $u$, the scanner satisfies
\[
R(A_tB_su)=B_tA_{2s+1}R(u),\qquad R(A_t)=B_t.
\tag{\ref{sec:a289128}.4}
\]
To verify the first equality, the initial $01$ pairs of $A_t$ emit
$(011)^t$ and leave one $0$. The first $0$ of $B_s$ combines with
that letter to form $00$, emitting $0010$. Each of the remaining
$s$ blocks $110$ is parsed as an unmatched $1$ followed by $10$,
and emits $1010$. Hence the output before $R(u)$ is
\[
(011)^t0010(10)^{2s}=0(110)^t0(10)^{2s+1}=B_tA_{2s+1}.
\]
The second equality follows from
$(011)^t0=0(110)^t$. These identities hold for arbitrary suffixes,
so they justify concatenating the particular blocks used below;
no concatenation property for arbitrary binary strings is assumed.

Put
\[
t_m=2^{m+1}-1,\qquad P_0=0010,\qquad
P_{m+1}=P_mA_{t_m}B_{t_m}.
\]
We prove simultaneously that
\[
R(P_mu)=R(P_m)R(u)\quad\hbox{for every }u,
\qquad R(P_m)=P_mA_{t_m}.
\tag{\ref{sec:a289128}.5}
\]
For $m=0$, the word $0010$ is consumed as the pairs $00,10$, and
$R(0010)=0010010=P_0A_1$. The scanner is at a complete boundary
after those input pairs. If (\ref{sec:a289128}.5) holds for $m$,
then (\ref{sec:a289128}.4) gives
\[
R(P_mA_{t_m}B_{t_m})
=P_mA_{t_m}B_{t_m}A_{2t_m+1}
=P_{m+1}A_{t_{m+1}},
\]
and the same computation with a suffix gives the boundary identity
at stage $m+1$. Therefore induction on the original iterates yields
\[
w_{2m+1}=P_m,\qquad w_{2m+2}=P_mA_{t_m}.
\tag{\ref{sec:a289128}.6}
\]
The next odd iterate appends $B_{t_m}$. Every word is thus a prefix
of its successor, including $w_0=00$.

\paragraph{Counting and the length recurrence.}
The block $A_t$ has length $2t+1$, with $t+1$ zeros and $t$ ones.
The block $B_t$ has length $3t+1$, with $t+1$ zeros and $2t$ ones.
Starting from $P_0=0010$, induction in
(\ref{sec:a289128}.6) proves all six formulas in
(\ref{sec:a289128}.1). In particular the lengths are unbounded,
so the nested prefixes determine $W$.

The length formulas for a fixed parity are a multiple of $2^m$
plus an affine function of $m$. They give the shorter relation
\[
L_{r+4}+2L_r=3L_{r+2}+3\qquad(r\ge1).
\]
Subtracting the relation at $r$ from the relation at $r+1$ gives
(\ref{sec:a289128}.2). Directly scanning the first five passes
gives its six initial values. This proves the conjectured equality
with the separately defined A289131 recurrence for every index.
Colin Barker already recorded the scalar parity formulas and
generating function in A289131 on July 2, 2017; the present argument
proves their equality with the original replacement lengths.

\paragraph{Infinite occurrence positions.}
Every positive-index iterate ends in $10$. Its final zero is
therefore the $Z_n$-th zero of the infinite word, and its final one
is the $O_n$-th one. Thus
\[
a(Z_n)=L_n,\qquad b(O_n)=L_n-1\qquad(n\ge1).
\tag{\ref{sec:a289128}.7}
\]
The count formulas show that the odd and even rank subsequences
increase strictly and tend to infinity. They also prove that both
symbols occur infinitely often. Substituting
(\ref{sec:a289128}.1) into (\ref{sec:a289128}.7) gives
\[
\begin{aligned}
\frac{a(Z_{2m+1})}{Z_{2m+1}}
 &=\frac{10\cdot2^m-3m-6}{4\cdot2^m-1},&
\frac{a(Z_{2m+2})}{Z_{2m+2}}
 &=\frac{14\cdot2^m-3m-7}{6\cdot2^m-1},\\
\frac{b(O_{2m+1})}{O_{2m+1}}
 &=\frac{10\cdot2^m-3m-7}{6\cdot2^m-3m-5},&
\frac{b(O_{2m+2})}{O_{2m+2}}
 &=\frac{14\cdot2^m-3m-8}{8\cdot2^m-3m-6}.
\end{aligned}
\]
Divide by $2^m$ and use $m/2^m\to0$ to obtain
(\ref{sec:a289128}.3). Were either full ratio convergent, both of
its displayed subsequences would have the same limit. Since
$5/2\ne7/3$ and $5/3\ne7/4$, neither converges. This refutes the
two proposed limits without relying on a finite-prefix
approximation. $\square$

\paragraph{Formal scope.}
Lean 4.19.0, with mathlib pinned at
\path{c44e0c8ee63ca166450922a373c7409c5d26b00b}, checks the entire
chain. The scanner uses actual lists of bits and recursively
consumes the original input. The exact block theorem is
\path{Kimberling289128.word_blocks}. Theorems
\path{word_nested}, \path{bit_stabilizes}, and
\path{infinite_word_unique} establish the nested limit. The
infinite letter at zero-based position $i$ is constructed as the
$i$-th letter of $P_i$; the proof gives $|P_i|\ge i+4$ and
stabilization in every iterate of index at least $2i+1$.

Both one-based position maps use \path{Nat.nth}. Their specification,
\path{position_spec}, proves the symbol and the exact preceding count
at every positive rank. Theorems \path{odd_positions} and
\path{even_positions} establish (\ref{sec:a289128}.7) for these maps.
The four real limits and strict growth of the rank sequences yield
\begin{center}
\path{position_ratios_do_not_converge}.
\end{center}
This theorem excludes every real limit for both ratios. The named
refutations are \path{A289129_limit_false} and
\path{A289130_limit_false}.

The length endpoint \path{original_length_conjecture} identifies
every literal length with the independent integer recurrence.
It includes the exceptional zeroth word and the exact start at
index six. All endpoint axiom dependencies are among
\path{propext}, \path{Classical.choice}, and \path{Quot.sound}.

\paragraph{Sources and prior-work screen.}
The package preserves the complete official OEIS exports for
A289128 revision 7, A289129 revision 4, A289130 revision 6, and
A289131 revision 19. Their saved Git blob hashes match the
official repository responses. The source review records searches
with both available engines, exact and semantic checks of the
current trureturing catalogue, and the existing mathematics
repository. No earlier proof of this exact length assertion or
refutation of these two exact position limits was located. This
bounded review does not certify historical priority.

The trureturing proof of A288476 uses the different rule
$00\to0101,1\to011$ and explicitly excludes the present position
and limiting-word questions. Michel Dekking's 2022 proof in
A288936 sends $10$ to $011$, rather than the $010$ used here, and
has different growth. The repository's earlier pair-rule and
morphic-discrepancy results also concern different substitutions.
Standard scanner recoding and subsequence methods are credited as
methods, not claimed as new general theory. The four limits above
are not asserted to classify all cluster points or determine exact
liminf and limsup values; corrected all-index inequalities are
outside this package's claimed scope.

\section{Verification, source correspondence, and attribution}
\label{sec:verification}

Each package uses Lean 4.19.0 and mathlib revision
\begin{center}
\texttt{c44e0c8ee63ca166450922a373c7409c5d26b00b}.
\end{center}
The repository contains source files, pinned Lake manifests,
complete mathematical READMEs, preserved primary-source exports,
bounded prior-publication reports, and separate author and
independent build records. There are fifteen local proof modules:
eight for A326926, one for A211417, and six for A289128 and its
companions. The default Lake targets include every local module.
From each package directory, the usual reproduction commands are
\begin{verbatim}
lake exe cache get
lake build
\end{verbatim}

Author and independent verification builds use fresh project
directories without copying any local proof objects. Only the
pinned external mathlib and dependency caches are reused. The
complete builds check the declared public endpoints from their
actual definitions. The formal files have no admissions, added
axioms, unsafe declarations, or native-decision proof shortcuts.
The reported endpoint dependencies are among the standard axioms
\path{propext}, \path{Classical.choice}, and \path{Quot.sound}.
This is kernel checking under the usual Lean trust boundary and
a separate computational review, not a claim of human peer review.

The most significant source-correspondence checks are:
\begin{itemize}
\item The polynomial package proves the derivative identity,
coprimality with the denominator, monic normalization, and exact
equality with the source's factorial-normalized reduced numerator.
Reversal is used only after proving the constant coefficient is
nonzero. The exceptional prime three and linear case are included.
\item The factorial-ratio package proves that natural division is
exact and equals the literal rational quotient. Its finite set
selects precisely the coprime shifts in the source, its constant
is explicitly positive, and its valuation argument controls the
whole product with multiplicities at every index.
\item The replacement package proves the literal scanner identity,
nesting and uniqueness of the infinite word, actual one-based
occurrence positions, and unbounded subsequence ranks. Its limits
are limits of the infinite position sequences. The length proof
retains all six initial values and begins the recurrence at six.
\end{itemize}

The seven preserved OEIS exports correspond to the exact entries
discussed above. Their revision timestamps and official Git blob
hashes are recorded in the source reports. The classification in
A326926 and A326933, and the length comment in A289128 and
A289131, are each counted once despite their repetition across entries. The three
values of $k$ in the factorial-ratio theorem are one stated general
family, not three separately counted resolutions.

The prior-work checks include both exact identifiers and semantic
searches using the relevant definitions. The polynomial review
credits Bala's source recurrence and Norfleet's classical
divisibility results, and compares related generalized Fibonacci
and tangent-polynomial irreducibility papers. The factorial-ratio
review excludes the known original integrality, previously solved
fixed factors, and the separately solved negative-shift family.
The word review credits Barker's scalar formulas and distinguishes
nearby replacement rules and earlier project packages.

No earlier proof of the precise five assertions claimed here was
located in the checked sources. That bounded search does not rule
out unindexed, unpublished, or differently formulated prior work.
No OEIS editorial acceptance, journal publication, or exhaustive
historical-priority conclusion is asserted. The mathematical
claims are the explicit theorems and refutations proved above.

\begin{thebibliography}{9}
\bibitem{A326926}
Clark Kimberling, OEIS A326926, introduced October 24, 2019;
Peter Bala's recurrence and divisibility addendum, November 8, 2026.
\url{https://oeis.org/A326926}.

\bibitem{A326933}
Clark Kimberling, OEIS A326933, introduced November 1, 2019.
\url{https://oeis.org/A326933}.

\bibitem{Norfleet}
Mark Norfleet, \emph{Characterization of Second-Order Strong
Divisibility Sequences of Polynomials}, Fibonacci Quarterly
43(2) (2005), 166--169.
\url{https://www.fq.math.ca/Papers1/43-2/paper43-2-12.pdf}.

\bibitem{A211417}
Peter Bala, OEIS A211417, introduced April 11, 2012;
positive-shift conjecture, August 28, 2025.
\url{https://oeis.org/A211417}.

\bibitem{A289128}
Clark Kimberling, OEIS A289128, June 28, 2017.
\url{https://oeis.org/A289128}.

\bibitem{A289129}
Clark Kimberling, OEIS A289129, June 28, 2017.
\url{https://oeis.org/A289129}.

\bibitem{A289130}
Clark Kimberling, OEIS A289130, June 28, 2017.
\url{https://oeis.org/A289130}.

\bibitem{A289131}
Clark Kimberling, OEIS A289131, June 28, 2017;
Colin Barker's scalar formulas, July 2, 2017.
\url{https://oeis.org/A289131}.
\end{thebibliography}
\end{document}
